\begin{tikzpicture}[font=\figurefont\small,text=BookInk,
  state/.style={draw=BookBlue,fill=BookBlue!5,minimum size=7mm,
    line width=.7pt,inner sep=1pt},
  action/.style={circle,draw=BookMuted,fill=white,minimum size=6mm,
    line width=.7pt,inner sep=1pt},
  edge/.style={-{Latex[length=1.8mm]},draw=BookMuted,line width=.8pt}]
  \node at (1.5,.6) {固定宽度};
  \node at (7,.6) {渐进加宽};
  \node[state] (fixed) at (1.5,-.1) {$s$};
  \node[action] (f1) at (.4,-1.5) {$a_1$};
  \node[action] (f2) at (2.6,-1.5) {$a_2$};
  \draw[edge] (fixed) -- (f1);
  \draw[edge] (fixed) -- (f2);
  \node[align=center,font=\figurefont\footnotesize,text=BookMuted] at (1.5,-2.5)
    {$N$增加，仍只查询\\预先保留的两个动作};
  \node[state] (wide) at (7,-.1) {$s$};
  \node[action] (w1) at (5.4,-1.5) {$a_1$};
  \node[action] (w2) at (7,-1.5) {$a_2$};
  \node[action,draw=BookTeal,fill=BookTeal!7] (w3) at (8.6,-1.5) {$a_3$};
  \draw[edge] (wide) -- (w1);
  \draw[edge] (wide) -- (w2);
  \draw[edge,BookTeal,dashed,line width=1.1pt] (wide) -- (w3);
  \node[align=center,font=\figurefont\footnotesize,text=BookMuted] at (7,-2.5)
    {教学规则 $|\mathcal A_N|\leq\lceil\sqrt N\rceil$\\
     $N:4\to9$，允许动作数 $2\to3$};
  \node[text=BookTeal,font=\figurefont\footnotesize] at (8.6,-1.99) {新增};
  \node[font=\figurefont\footnotesize,text=BookMuted] at (4.7,-3.4)
    {加宽动作需要覆盖条件；随机后继的加宽是另一层操作};
\end{tikzpicture}
